\subsection{MIDI}\label{subsec:midi}

\textbf{MIDI} (Musical Instrument Digital Interface) is a standardized protocol for transmitting musical
control and timing information between electronic instruments and computers in real time~\cite{Moog_1986}.
Rather than encoding audio signals, MIDI represents music as a stream of discrete digital messages describing
performance actions such as note onset, pitch, and control changes.

%% ---------------------------------------------------------------------------------------------------------------------

\subsubsection{Key Features}

MIDI encodes musical information as structured messages consisting of status and data bytes~\cite{Moog_1986}.
A fundamental example is the \textit{Note-On} message, which specifies both a pitch (key number) and a
velocity value corresponding to how quickly a key is pressed~\cite{Moog_1986}.
These values are discretized into integer ranges (0–127), reflecting MIDI’s 7-bit resolution~\cite{Moog_1986}.

Additional message types include \textit{Note-Off}, \textit{Control Change}, \textit{Program Change}, and
\textit{Pitch Bend}, each representing different aspects of performance~\cite{Moog_1986}.
Messages are transmitted serially at a fixed rate of 31.25 kbit/s, which constrains the amount of data that
can be sent in real time~\cite{Moog_1986}.

Furthermore, MIDI organizes communication into 16 logical channels, allowing multiple instruments to be
controlled over a single connection~\cite{Moog_1986}.

%% ---------------------------------------------------------------------------------------------------------------------

\subsubsection{Limitations}

For the purposes of this thesis, the main limitation of MIDI is not utility, but scope.
Its design is rooted in a keyboard-centric model of musical interaction, where music is represented as
discrete note events supplemented by parameterized controls such as velocity and aftertouch~\cite{Moore_1988}.
This makes it difficult to represent continuous musical gestures without approximating them as streams of
discrete messages~\cite{Moore_1988}.

Moore (1988) argues that MIDI conflates musical structure with performance data and lacks the ability to
represent higher-level abstractions such as phrasing, hierarchy, and harmonic relationships~\cite{Moore_1988}.
As a result, MIDI functions primarily as a performance protocol rather than a compositional
representation~\cite{Moore_1988}.

Additionally, the finite resolution of MIDI data (typically 7-bit, or 14-bit for certain controls) introduces
quantization effects, particularly noticeable in pitch and dynamics~\cite{Moore_1988}.

Timing is another critical limitation: because MIDI messages are transmitted serially, dense musical passages
can introduce latency and jitter~\cite{Moog_1986}.
Even under typical conditions, delays on the order of 10–30 milliseconds may occur due to processing and
transmission constraints~\cite{Moog_1986}.

%% ---------------------------------------------------------------------------------------------------------------------

\subsubsection{Relevance to Our DSL}

MIDI highlights the strengths and weaknesses of low-level, event-based representations of music.
While it excels at interoperability and real-time performance control, it lacks the expressive power needed
to encode higher-level musical structures.

This distinction is important for the architecture proposed in this thesis.
MIDI remains a valuable backend target for rendering and interchange, but it is not sufficient as the primary
representation for structured composition.
This motivates a DSL that separates compositional intent from performance encoding, provides richer abstractions for
musical structure, and can ultimately compile to formats such as MIDI for execution.
